\scanpage{3}\mspage{1}
% archive sheet number: 2 (upper left)
% l. 1
\uline{Platitude d'une adhérence schématique}
\margin{$U\cap T\;\unsure{\frac{\phantom{T}}{|X|}}$}

% ll. 2–3. The i is beneath the inclusion sign.
1.1.
\diagram{\begin{tikzcd}[ampersand replacement=\&,column sep=small]
U\mathrel{\underset{i}{\subset}} \&[-1.5em] X\arrow[d,"f"] \& F\\
\& Y \&
\end{tikzcd}}
données $f\colon\inserted{X\to Y}$, $U$ ouvert \inserted{rétrocompact} de $X$, $F$ quasi-cohérent sur $X$
$G_U$ quotient de $F|U$, quasi-cohérent et \uline{plat} sur $Y$

% ll. 4–9
\uline{Pb} Trouver un quotient \inserted{\unsure{qu. coh.}} $G$ de $F$, qui \uline{prolonge}
$G_U$, \deleted{\unsure{et}} qui soit plat sur $Y$, et \deleted{tel} que
\unsure{i.e.} $G\xrightarrow{\inserted{\text{rel. à }i}}G_U$ \inserted{$[\,G\to i_*(G_U)$} soit \uline{univ. séparant} rel $/S$.
[Il \inserted{faut et} suffit pour
ceci, \deleted{\unsure{lorsque} $G$} si $f$ est loc. de présentation finie, et si
\deleted{$Z=\operatorname{Supp}G$, que} $G$ est de prés. finie sur $S$, que
pour tout $y\in Y$, $U_y$ contienne $\operatorname{Ass}G_y$].

% margin note written sideways
\margin{\deleted{On peut donc supposer $X$ plat sur $S$ (lisse sur $S$ si à l'origine il était \inserted{loc.} de type fini) et $F$ libre sur $X$ \deleted{\unsure{coh.}} (et de type fini s'il l'était initialement), quitte à \unsure{coiffer|coffer}}}

% ll. 10–15. The lower arrows belong to the image formula.
Si le Pb a une solution, elle est unique,
savoir
\diagram{\begin{tikzcd}[ampersand replacement=\&,column sep=small,row sep=small]
G\simeq\operatorname{Im}\bigl( F \arrow[r] \arrow[dr] \& i_*(G_U)\bigr).\\
\& i_*(F_U)\arrow[u]
\end{tikzcd}}
La question est donc de savoir si le $G$ défini par cette formule
est \deleted{\unsure{univ.}}\ \deleted{\unsure{univ.}} plat rel $/Y$ et $G\to G_U$ est
univ. \deleted{\unsure{injec.}} séparant rel $/Y$. \inserted{\uline{Remarque} Question \uline{locale} sur $X$ et $Y$, qu'on peut supposer affines !}
L'unicité implique
aussi

% ll. 16–18
\uline{Lemme} 1.2. Question locale \inserted{\uline{sur $Y$}} pour fpqc. De \unsure{m} la
\margin{SGA 1 \uline{VIII}}
question analogue, où on veut que $G$ soit \deleted{loc} de
prés. finie sur $X$. // Plus gén., descente possible \uline{sur $X$} pour morphisme fid. plat \unsure{qu. cpt|qu. cpct} $X'\to X$.

% ll. 19–20
La descente non plate \inserted{(du quotient \uline{plat} d'un faisceau déjà donné),} permet ensuite de prouver :

% ll. 21–23
\uline{Lemme} 1.3. Supposons $Y$ artinien \inserted{$=\operatorname{Spec}A$}, \deleted{\unsure{q}} et qu'il existe un
épimorphisme \inserted{$A\hookrightarrow\Gamma(Y,\mathcal O_Y)$} de schémas $Y'\to Y$, tel que après ce
changement de base, le Pb ait une solution. Supposons
% slanted
\margin{Il faut supposer soit $Y'\to Y$ fini ou $f\colon X\to Y$ loc. de \unsure{p.f.|prés. f.}, et on cherche des solutions de prés. finie}
